%% =====================================================================
\section{Experimental Methodology}
\label{sec:method}
%% =====================================================================

\subsection{Environment}
All experiments ran on a single Windows~11 Pro workstation under Docker Desktop
(Docker Engine 29.8.2), with the nine containers of Table~\ref{tab:services} on one
Compose bridge network. Container resource limits were left at Docker Desktop defaults.
Running every component on one host removes network variability between tiers and
makes runs comparable, at the cost of shared CPU contention between the attacker,
the victim, and the instrumentation (Section~\ref{sec:validity}).

\subsection{Experimental Conditions and Protocol}
Experiments follow a three-condition protocol, executed from the dashboard presets:
\begin{description}[leftmargin=1.2em,style=nextline]
\item[E1 -- Baseline.] Normal traffic only, defenses off.
\item[E2 -- Attack, undefended.] Normal traffic plus the rotating burst attack, both
defenses off.
\item[E3 -- Attack, defended.] Normal traffic plus the attack, with cache locking,
application rate limiting, and edge rate limiting enabled.
\end{description}
A typical session cycles E1~$\rightarrow$~E2~$\rightarrow$~E1~$\rightarrow$~E3 with
roughly 5--20~s per condition, so that each undefended and each defended attack is
preceded by a fresh normal period from which the detectors can (re)learn the
baseline. In addition, the final session contains one long defended run spanning
${\approx}8.1$~h of wall-clock time, used as a soak test of the defense; it contains
4.5~h of valid samples, the remainder being 15 gaps during which the host was asleep.
Two further procedures were used:
\begin{itemize}[leftmargin=*]
\item \textbf{Ablation.} Starting from E2, defenses were enabled one at a time (cache
locking only; plus stale-while-revalidate; plus edge rate limiting) while observing the
dashboard and the CacheShop storefront (Section~\ref{sec:ablation}). Because the
defense implementation evolved during the project, the 14 sessions that include
defended traffic also form a natural longitudinal ablation, analyzed in
Section~\ref{sec:longitudinal}.
\item \textbf{Automated validation.} A 16-case end-to-end test suite drives the system
through the same control API and checks functional and quantitative success criteria
(Section~\ref{sec:validation}).
\end{itemize}

\subsection{Dataset}
The monitor writes one CSV row per sampling interval (Table~\ref{tab:csv}). A new file
is started each time the monitor container boots. The study uses all 18 files produced
on 22~October~2026, totalling 92{,}531 rows and ${\approx}28$~h of wall-clock logging,
much of it idle (Table~\ref{tab:dataset}). The sampling interval was 1~s in the first
eight sessions and 0.5~s thereafter.

\begin{table}[t]
\centering
\caption{Columns of the monitor's CSV log (one row per sampling interval).}
\label{tab:csv}
\small
\begin{tabularx}{\linewidth}{@{}L{2.6cm}Y@{}}
\toprule
Column & Definition \\
\midrule
\texttt{ts} & Unix timestamp of the sample \\
\texttt{mode} & idle / normal / attack / protected (attack with $\ge$1 defense on) \\
\texttt{rps} & Application requests per second (all outcomes, incl.\ 429 and 503) \\
\texttt{db\_qps} & Database query \emph{attempts} per second \\
\texttt{cache\_hits}, \texttt{cache\_misses} & Counts in the interval \\
\texttt{hit\_rate} & hits / (hits + misses) in the interval \\
\texttt{resp\_avg\_ms} & Mean in-application latency of requests completed in the interval \\
\texttt{resp\_p95\_ms} & Histogram-bucket approximation of p95 latency \\
\texttt{failed} & \emph{Count} of HTTP~503 responses in the interval (not a rate) \\
\texttt{rule\_alert}, \texttt{ml\_alert} & Detector outputs (0/1) \\
\texttt{ml\_score} & Normalized anomaly score for display (0.08 = model not yet trained) \\
\bottomrule
\end{tabularx}
\end{table}

\begin{table}[t]
\centering
\caption{Dataset summary. ``Final'' is the session with the complete defense
(edge rate limiting active), on which the main results are computed.}
\label{tab:dataset}
\small
\begin{tabular}{@{}lrrrrr@{}}
\toprule
Subset & Rows & Idle & Normal & Attack & Defended \\
\midrule
All 18 sessions & 92{,}531 & 50{,}732 & 1{,}477 & 1{,}648 & 38{,}674 \\
Final session (16:58), valid rows & 61{,}923 & 25{,}586 & 332 & 733 & 35{,}272 \\
\bottomrule
\end{tabular}
\end{table}

\paragraph{Choice of the primary dataset.} The defense evolved during the day
(Section~\ref{sec:lessons}). Only in the last session is the edge rate limiter
demonstrably active: the application's request rate in defended mode is pinned at the
configured 50~req/s (median 50, IQR 47--53), whereas in every earlier session it is
170--350~req/s. We therefore report the main results on the final session (12
undefended attack episodes totalling 13~min of valid samples, 10 defended episodes
totalling 5.6~h of which 4.5~h are the soak run, and 12 normal episodes) and use the earlier sessions for the longitudinal
analysis and for robustness checks.

\subsection{Pre-processing}
\begin{enumerate}[leftmargin=*]
\item \textbf{Interval length.} For each row, $\Delta t$ is the difference between
consecutive timestamps. Rows with $\Delta t\notin(0.25, 2.0)$~s are excluded (245 of
92{,}531 rows, 0.26\%). These arise when the host sleeps or the monitor stalls and produce
meaningless rates; for example one such row in the soak run reports a mean latency of
over 1{,}000~s.
\item \textbf{Failure rate.} Because \texttt{failed} is a per-interval count, the
failure rate is computed as $\texttt{failed}/\Delta t$.
\item \textbf{Miss ratio.} Computed as misses/(hits+misses) per interval, 0 when both
are zero.
\item \textbf{Episodes.} Each maximal run of consecutive rows with the same mode is an
\emph{episode}. Episodes shorter than 8 rows are excluded from per-episode analyses
(they are transitions caused by rapid button presses).
\item \textbf{ML readiness.} Rows with \texttt{ml\_score}$=0.08$ and no ML alert are
those in which the forest had not yet been trained; per-episode ML results report
whether the model was ready.
\end{enumerate}

\subsection{Ground Truth and Evaluation Metrics}
The mode recorded in each row provides ground truth: intervals in \emph{attack} mode
are positives (the attack is running and undefended) and intervals in \emph{normal}
mode are negatives. Defended and idle intervals are reported separately, because the
correct answer there is a matter of definition: during a defended attack the attack
\emph{is} present but causes no harm, and during idle periods there is no traffic at
all. Table~\ref{tab:metric-defs} defines the evaluation metrics.

\begin{table}[t]
\centering
\caption{Evaluation metrics.}
\label{tab:metric-defs}
\small
\begin{tabularx}{\linewidth}{@{}L{3.6cm}Y@{}}
\toprule
Metric & Definition \\
\midrule
DB load & Median (IQR, p95, max) of \texttt{db\_qps} over intervals in a condition \\
Latency & Median (IQR, p95, max) of per-interval mean latency \\
Failure rate & HTTP 503 per second; fraction of intervals with any 503 \\
Reduction factor & Ratio of medians (undefended / defended), 95\% bootstrap CI \cite{efron1994} \\
Coalescing ratio $\rho$ & Total cache misses / total DB attempts in a condition \\
Completed DB queries & $\max(0, q_{db}\Delta t - n_{503})/\Delta t$ per interval \\
Detection latency & Time from first interval of an attack episode to first alert \\
Episode recall & Fraction of attack episodes with at least one alert \\
Interval precision / recall / F1 & Over attack (positive) and normal (negative) intervals \\
False-positive rate & Alerts in normal intervals / normal intervals \\
ROC AUC & Area under the ROC curve of a continuous score \cite{fawcett2006} \\
Recovery time & First interval of a new condition to the first of two consecutive healthy intervals (Eq.~\ref{eq:healthy}) \\
\bottomrule
\end{tabularx}
\end{table}

\subsection{Offline Analyses}
Three analyses go beyond the detectors' live output.
\begin{itemize}[leftmargin=*]
\item \textbf{Threshold sensitivity of the rule.} Equation~\eqref{eq:rule} is
re-evaluated on the logged features over a grid of $K\in\{2,3,5,8,12\}$,
$\tau_L\in\{100,150,200,300\}$~ms and $\tau_m\in\{0.3,0.5,0.7\}$, using the median
normal database rate of the final session (11~q/s) as baseline.
\item \textbf{Single-signal baselines.} Each condition of the rule, a ``failures
present'' detector, and a pure volume threshold (rps$\,>100$) are evaluated alone, both
on the attack/normal intervals and on 2{,}080 \emph{absorbed-burst} intervals: defended
intervals from six earlier sessions in which the full attack volume reached the
application (median 280~req/s) but was served from cache with low database load
(median 2~q/s) and low latency (median 26~ms). These are high-volume, harmless
intervals---a laboratory proxy for a flash crowd that the cache absorbs.
\item \textbf{Offline IsolationForest replay.} The detector pipeline is re-run offline
with the same features, scaler, and hyper-parameters: it is trained on the first half
of the final session's normal intervals (166) and tested on the second half (166) plus
all 733 attack intervals. This permits ROC analysis, a sweep of the
\texttt{contamination} parameter, a sweep of the training-window size, and a
random-seed stability check.
\end{itemize}
All analysis code uses pandas, NumPy, scikit-learn, and Matplotlib, and reads only the
CSV logs.
